\documentclass[conference]{IEEEtran}

\usepackage[T1]{fontenc}
\usepackage{amsmath,amssymb}
\usepackage{booktabs}
\usepackage{graphicx}
\usepackage{float}
\usepackage{xcolor}
\usepackage{listings}
\usepackage{enumitem}
\usepackage{hyperref}

% Configure listings for Python Code
\lstdefinestyle{pythonstyle}{
    language=Python,
    basicstyle=\ttfamily\scriptsize,
    columns=fullflexible,
    keepspaces=true,
    breaklines=true,
    breakatwhitespace=false,
    frame=single,
    framerule=0.3pt,
    rulecolor=\color{black!30},
    numbers=left,
    numberstyle=\tiny\color{gray},
    numbersep=4pt,
    xleftmargin=14pt,
    framexleftmargin=10pt,
    aboveskip=4pt,
    belowskip=4pt,
    showstringspaces=false,
    captionpos=b
}

% Configure listings for Console Outputs / Execution Traces
\lstdefinestyle{outputstyle}{
    language={},
    basicstyle=\ttfamily\tiny,
    columns=fullflexible,
    keepspaces=true,
    breaklines=true,
    breakatwhitespace=false,
    frame=single,
    framerule=0.3pt,
    rulecolor=\color{black!25},
    backgroundcolor=\color{black!3},
    numbers=none,
    xleftmargin=6pt,
    framexleftmargin=6pt,
    aboveskip=3pt,
    belowskip=5pt,
    showstringspaces=false,
    captionpos=b
}

\lstset{style=pythonstyle}

% Layout and Spacing Settings
\setlength{\textfloatsep}{5pt plus 1pt minus 2pt}
\setlength{\floatsep}{4pt plus 1pt minus 2pt}
\setlength{\intextsep}{4pt plus 1pt minus 2pt}
\setlength{\abovecaptionskip}{3pt}
\setlength{\belowcaptionskip}{0pt}
\setlength{\tabcolsep}{3.2pt}
\renewcommand{\arraystretch}{0.92}
\setlength{\parskip}{0pt}
\setlength{\partopsep}{0pt}

\setlist{nosep,leftmargin=*}

\hypersetup{
    colorlinks=true,
    linkcolor=blue!70!black,
    urlcolor=blue!70!black,
    citecolor=blue!70!black
}

\title{%
\textbf{Build for Bharat 2.0 Hackathon:}\\
\textit{Approach Notice Submission}\\[0.45em]
\large\textbf{Team: Command Q}\\
\normalsize Team Leader: Rajat Kumar Gupta (25BDA70384)\\[0.7em]
\large\textbf{Career Navigation System for Data Talent Workforce:}\\
\large A Data-Driven Approach to Skills Alignment and Career Progression
}

\author{\IEEEauthorblockN{Command Q}}

\begin{document}

\maketitle

\section{Problem Definition and Analytics Objective}

\subsection{Problem Identification}

The data science and analytics workforce in India and globally faces a critical structural challenge: aspiring and practicing professionals struggle to navigate career progression effectively due to profound information asymmetry. Despite surging demand for data talent, individuals lack empirical clarity regarding:
\begin{itemize}
    \item \textbf{Personality-Role Congruence}: How foundational psychometric traits dictate long-term success, role sustainability, and progression in high-ambiguity senior technical leadership.
    \item \textbf{Skills Valuation Disconnect}: Which specific competencies drive internal promotion velocity versus advertised hiring demand, avoiding market signaling distortions.
    \item \textbf{Strategic Human Capital Investment}: How to allocate finite reskilling hours to maximize competitive differentiation and wage appreciation.
    \item \textbf{Employer-Experience Alignment}: Which organizations offer optimal compensation brackets and liquidity for an individual's specific career tenure.
\end{itemize}

This pervasive information asymmetry precipitates career churn, misdirected upskilling expenditure, and wage suppression. Our study provides a rigorous, empirical foundation for an intelligent career navigation system tailored to three primary career cohorts: Senior Data Scientists (SDS), Junior Data Scientists (JDS), and Analytics/BI Professionals.

\subsection{Scope and Analytical Depth}

Our analytical architecture integrates four distinct empirical dimensions:
\begin{enumerate}
    \item \textbf{Psychometric-Career Alignment}: Quantifying Big Five personality trait effect sizes ($d > 1.70$) and predictive classification bounds for senior data science leadership.
    \item \textbf{Skills Market Intelligence}: Evaluating skill-outcome effect sizes, market penetration rates across 14,749 postings, and econometric wage premiums.
    \item \textbf{Strategic Market Positioning}: Establishing a $2 \times 2$ Strategic Skills Matrix and Rank Gap metric to identify undervalued high-leverage skills.
    \item \textbf{Econometric Wage Modeling \& Matching}: Developing cross-validated salary regression models ($R^2 = 0.571$) and a multi-attribute organizational matching engine.
\end{enumerate}

\subsection{Analytics Objective and Success Criteria}

The primary objective is to build an inspectable, data-driven career navigation engine that operationalizes these findings into actionable career prescriptions.
\textbf{Measurable Evaluation Criteria}:
\begin{itemize}
    \item Statistical significance and large effect sizes ($|d| \ge 0.80$, AUC $> 0.75$, $p < 0.001$) across validated traits and skills.
    \item Cross-validated classification accuracy ($> 85\%$ ROC-AUC) demonstrating high discriminative validity.
    \item High goodness-of-fit ($R^2 > 0.50$, MAE reduction $> 35\%$) in baseline salary modeling.
    \item Total reproducibility where every metric, decision tree rule, and regression coefficient is traced to open code listings and verified console outputs.
\end{itemize}

\section{Approach Description and Prototype Architecture}

\subsection{Overall System Architecture}

Our proposed prototype---the \textbf{Career Navigation System for Data Talent Workforce}---translates empirical findings into a four-stage interactive career navigation pipeline, visualized in Fig.~\ref{fig:prototype_arch}.

\begin{figure*}[!t]
\centering
\includegraphics[width=0.95\textwidth]{figures/prototype_architecture.png}
\caption{Career Navigation System: End-to-End Modular Prototype Architecture.}
\label{fig:prototype_arch}
\end{figure*}

The architecture couples user assessment inputs with four specialized analytical engines:
\begin{enumerate}
    \item \textbf{Career Track Norming \& Psychometric Fit Engine}: Standardizes candidate profile inputs against benchmark cohorts.
    \item \textbf{Strategic Skills Recommendation Engine}: Evaluates proficiency against the $2 \times 2$ Strategic Market Matrix.
    \item \textbf{Progression Path Optimizer}: Traverses empirical skill co-development graphs to design personalized learning sequences.
    \item \textbf{Multi-Attribute Company Matching Engine}: Matches candidates to hiring organizations optimizing experience fit, salary bracket, and vacancy liquidity.
\end{enumerate}

\subsection{Mathematical Formulations of Recommendation Engines}

\subsubsection{Engine 1: Psychometric Fit Engine}
To assess alignment with senior data science leadership without arbitrary heuristics, we formulate a normalized psychometric fit index ($S_{\text{fit}} \in [0, 100]$):
\begin{equation}
S_{\text{fit}} = 100 \times \sum_{i=1}^{5} w_i \cdot \max\left(0, \min\left(1, \frac{x_i - \mu_{\text{low}, i}}{\mu_{\text{high}, i} - \mu_{\text{low}, i}}\right)\right)
\label{eq:fit_score}
\end{equation}
where trait weights $w_i$ are proportional to empirical effect sizes (Cohen's $d_i$):
\begin{equation}
w_i = \frac{|d_i|}{\sum_{j=1}^5 |d_j|}
\end{equation}
Empirically calibrated weights from our benchmark cohort ($N=161$) yield:
\begin{itemize}
    \item Conscientiousness: $w_C = 0.30$ ($\mu_{\text{high}} = 53.68, \mu_{\text{low}} = 35.74$)
    \item Openness to Experience: $w_O = 0.29$ ($\mu_{\text{high}} = 48.49, \mu_{\text{low}} = 33.32$)
    \item Extraversion: $w_E = 0.19$ ($\mu_{\text{high}} = 48.86, \mu_{\text{low}} = 36.88$)
    \item Agreeableness: $w_A = 0.10$ ($\mu_{\text{high}} = 47.72, \mu_{\text{low}} = 41.12$)
    \item Neuroticism: $w_N = 0.12$ ($\text{Penalizes extreme scores} > 50$)
\end{itemize}

\subsubsection{Engine 2: Strategic Skills Recommendation Engine}
To balance internal promotion return against external market employability, the system computes a Skill Priority Score $P_k$ for candidate skill deficits:
\begin{equation}
P_k = \alpha \cdot d_k + \beta \cdot (1 - D_k) + \gamma \cdot \text{RankGap}_k
\label{eq:skill_priority}
\end{equation}
where $d_k$ is the standardized outcome effect size, $D_k \in [0, 1]$ is the macro market demand rate, and $\text{RankGap}_k = \text{DemandRank}_k - \text{EffectRank}_k$. Calibrated hyperparameters ($\alpha=0.45, \beta=0.30, \gamma=0.25$) prioritize "Hidden Gem" competencies that unlock disproportionate promotion velocity.

\subsubsection{Engine 3: Progression Path Optimizer}
Reskilling paths are optimized via a directed acyclic graph $G=(V, E)$, where vertices represent skills and edge weights represent empirical co-development correlations $r_{ij}$:
\begin{equation}
W(e_{ij}) = r_{ij} \cdot \left(1 + \mathbb{I}(\text{Exp} \ge 4) \cdot \Delta D_{\text{senior}, j}\right)
\label{eq:graph_weight}
\end{equation}
This accounts for the structural demand shift where Big Data infrastructure skills become mandatory only beyond Year 4.

\subsubsection{Engine 4: Multi-Attribute Company Matching Engine}
For candidate experience $E$, the system scores prospective hiring organizations $c$:
\begin{equation}
M_c = 0.50 \cdot \text{ExpFit}(E, c) + 0.35 \cdot \text{SalaryScore}(c) + 0.15 \cdot \text{LiquidityScore}(c)
\label{eq:match_score}
\end{equation}
where:
\begin{align}
\text{ExpFit}(E, c) &= \begin{cases} 
1.0 & \text{if } \text{MinExp}_c \le E \le \text{MinExp}_c + 2 \\ 
\max(0, 1 - 0.25 \cdot |\text{MinExp}_c - E|) & \text{otherwise} 
\end{cases} \\
\text{SalaryScore}(c) &= \text{Percentile}\left(\text{AvgSalary}_c \mid \text{Seniority Tier}\right) \\
\text{LiquidityScore}(c) &= \frac{\ln(1 + \text{Openings}_c)}{\max_j \ln(1 + \text{Openings}_j)}
\end{align}

\subsection{User Personas and Journey Flows}

The prototype supports three archetypal user journeys:
\begin{itemize}
    \item \textbf{Persona 1: Junior Aspiring Data Scientist (1--2 YOE)}: Assesses foundational coding/ML skills. Identifies Storytelling deficit ($\text{RankGap}=+2$). Prescribes data narrative modules; routes toward high-hike organizations.
    \item \textbf{Persona 2: Mid-Career Business/Data Analyst (3--5 YOE)}: Tests transition potential to Data Science. Flags math/statistics gap and identifies Big Data timing inflection. Targets firms paying $+72\%$ senior analyst premiums.
    \item \textbf{Persona 3: Senior Technical Specialist (7+ YOE)}: Evaluates transition into Lead Data Scientist or Data Architect roles. Psychometric engine assesses Conscientiousness and Openness, identifying leadership and architectural readiness.
\end{itemize}

\section{Data Exploration and Preprocessing}
\phantomsection\label{sec:data_exploration}

\subsection{Data Sources Overview}

Our empirical investigation utilizes four core datasets comprising 16,741 raw records:
\begin{table}[!t]
\centering
\caption{Summary of Empirical Datasets}
\label{tab:datasets}
\footnotesize\begin{tabular}{@{}llrr@{}}
\toprule
\textbf{Dataset} & \textbf{Core Construct} & \textbf{Raw Rows} & \textbf{Core Rows} \\
\midrule
SDS Personality & Big Five traits (OCEAN), Success Flag & 161 & 161 \\
JDS Skills & 5 Skill Proficiencies, Salary Hike Flag & 139 & 139 \\
Data Science Jobs & Experience, Salary LPA, Openings & 1,602 & 1,595 \\
Analytics Jobs & Designations, Skills, Experience, Salary & 14,839 & 14,749 \\
\bottomrule
\end{tabular}
\end{table}

\subsection{Data Quality Assessment and Treatment Policy}
\phantomsection\label{sec:data_quality}\phantomsection\label{sec:treatment}

Prior to modeling, we conducted systematic audits across schema completeness, key integrity, duplicate rows, and domain constraints \hyperref[lst:audit]{[Listing G.3 \& Output G.3]}. Crucially, our analysis verified that ID columns across datasets represent coincidental numerical overlaps rather than foreign keys \hyperref[lst:domain]{[Listing G.4 \& Output G.4]}. Analysis therefore proceeds at the construct level.

\begin{table}[!t]
\centering
\caption{Data Cleaning and Treatment Policy}
\label{tab:cleaning_policy}
\footnotesize\begin{tabular}{@{}lrrrl@{}}
\toprule
\textbf{Source} & \textbf{Raw} & \textbf{Clean} & \textbf{Dropped} & \textbf{Rationale} \\
\midrule
JDS Skills & 139 & 139 & 0 (0.0\%) & Fully valid scores (1--5) \\
SDS Personality & 161 & 161 & 0 (0.0\%) & Fully valid traits (17--68) \\
Data Science Jobs & 1,602 & 1,595 & 7 (0.4\%) & Placeholder company flags \\
Analytics Jobs & 14,839 & 14,749 & 90 (0.6\%) & Gig/data-entry deduplication \\
\bottomrule
\end{tabular}
\end{table}

As documented in Table~\ref{tab:cleaning_policy} and verified in \hyperref[lst:treatment]{[Listing G.5 \& Output G.5]}, we applied conservative treatments that retained $>99.4\%$ of records while purging corrupt records and gig postings.

\subsection{Distributional Characteristics and Normality Testing}
\phantomsection\label{sec:normality}

We conducted Shapiro-Wilk normality tests on all numerical attributes \hyperref[lst:normality]{[Listing G.7 \& Output G.7]}. Every single trait and skill score rejected normality at $p < 0.001$ (e.g., JDS Coding $W=0.788, p=1.2\times 10^{-13}$; SDS Conscientiousness $W=0.921, p=8.9\times 10^{-8}$). Consequently, we anchored all subsequent hypothesis testing in non-parametric Mann-Whitney $U$ formulations.

The empirical distributions are presented across Figs.~\ref{fig:sds_traits}, \ref{fig:jds_skills}, \ref{fig:dsj_dist}, \ref{fig:analytics_dist}, and \ref{fig:categorical}.

\begin{figure*}[!t]
\centering
\includegraphics[width=\textwidth]{figures/sds_trait_distributions.png}
\caption{Distribution of Big Five Personality Traits for Senior Data Scientists.}
\label{fig:sds_traits}
\end{figure*}

\begin{figure*}[!t]
\centering
\includegraphics[width=\textwidth]{figures/jds_skill_distributions.png}
\caption{Proficiency Distributions Across Five Core Junior Data Science Skills.}
\label{fig:jds_skills}
\end{figure*}

\begin{figure}[!t]
\centering
\includegraphics[width=\columnwidth]{figures/datascience_jobs_distributions.png}
\caption{Data Science Jobs: Salary LPA and Experience Distributions.}
\label{fig:dsj_dist}
\end{figure}

\begin{figure*}[!t]
\centering
\includegraphics[width=\textwidth]{figures/analytics_jobs_distributions.png}
\caption{Analytics Jobs: Market Experience and Salary Overview.}
\label{fig:analytics_dist}
\end{figure*}

\begin{figure}[!t]
\centering
\includegraphics[width=\columnwidth]{figures/categorical_overview.png}
\caption{Categorical Feature Distributions Across Job Market Postings.}
\label{fig:categorical}
\end{figure}

\subsection{Feature Derivation and Construct Grounding}
\phantomsection\label{sec:data_derivation}

We engineered standardized analytical constructs \hyperref[lst:features]{[Listing G.6 \& Output G.6]}:
\begin{itemize}
    \item \textbf{Role Family Categorization}: Clustered unstructured job designations into 6 core tracks: Analyst ($n=4,007$), Consultant/Manager ($n=3,726$), Data Engineer ($n=446$), Data Scientist ($n=439$), Data Architect ($n=357$), and AI/ML Engineer ($n=201$).
    \item \textbf{Skill Dimension Extraction}: Extracted standardized keyword vectors for Big Data, Math/Stats, Coding, AI/ML, and Dashboard/Storytelling.
    \item \textbf{Seniority Standardization}: Mapped experience into Junior (0--3 yrs), Mid (3--7 yrs), and Senior (8+ yrs).
\end{itemize}

\section{Implementation Pipeline and Reproducibility}
\phantomsection\label{sec:pipeline}

The entire analytical workflow is organized into an inspectable, staged Python architecture \hyperref[lst:setup]{[Listing G.1]}, separating data ingestion \hyperref[lst:loading]{[Listing G.2 \& Output G.2]}, auditing, transformation, statistical testing, multivariate modeling, and artifact export \hyperref[lst:export]{[Listing G.13 \& Output G.13]}.

All derived datasets, bridge tables, and high-resolution figures are automatically generated and validated via the final recap engine \hyperref[lst:recap]{[Listing G.14 \& Output G.14]}, ensuring full reproducibility across research and enterprise deployments.

\section{Data Analysis and Empirical Foundations}

\subsection{Personality-Career Alignment Analysis}
\phantomsection\label{sec:personality_analysis}\phantomsection\label{sec:sds_models}

\subsubsection{Non-Parametric Group Differences}
Evaluating Big Five personality scores between high-performing senior data scientists ($n=85$) and the comparison cohort ($n=76$) reveals massive, statistically robust effect sizes (Table~\ref{tab:sds_effects}, verified in \hyperref[lst:sds_effects]{[Listing G.8 \& Output G.8a]}).

\begin{table}[!t]
\centering
\caption{Personality Trait Effect Sizes for Senior Data Scientists}
\label{tab:sds_effects}
\footnotesize\begin{tabular}{@{}lrrrrc@{}}
\toprule
\textbf{Trait} & \textbf{High $\mu$} & \textbf{Low $\mu$} & \textbf{Cohen's d} & \textbf{AUC} & \textbf{Bonferroni $p$} \\
\midrule
Conscientiousness & 53.68 & 35.74 & 1.815 & 0.878 & $1.44 \times 10^{-16}$* \\
Openness to Exp. & 48.49 & 33.32 & 1.774 & 0.885 & $3.64 \times 10^{-17}$* \\
Extraversion & 48.86 & 36.88 & 1.115 & 0.783 & $5.79 \times 10^{-10}$* \\
Agreeableness & 47.72 & 41.12 & 0.595 & 0.652 & 0.004* \\
Neuroticism & 36.13 & 36.26 & -0.012 & 0.534 & 1.000 (ns) \\
\bottomrule
\end{tabular}
\end{table}

\begin{figure*}[!t]
\centering
\includegraphics[width=\textwidth]{figures/sds_outcome_comparison.png}
\caption{Personality Trait Differences Between High and Low Performing Senior Data Scientists.}
\label{fig:sds_outcome}
\end{figure*}

\begin{figure}[!t]
\centering
\includegraphics[width=\columnwidth]{figures/sds_correlation.png}
\caption{Inter-Trait Correlation Matrix for Senior Data Scientists.}
\label{fig:sds_corr}
\end{figure}

Conscientiousness ($d=1.815, \text{AUC}=0.878$) and Openness to Experience ($d=1.774, \text{AUC}=0.885$) demonstrate extraordinary discriminatory power. Senior technical leadership requires intellectual curiosity to formulate non-standard algorithmic approaches coupled with rigorous conscientiousness to enforce pipeline reliability and business delivery.

\subsubsection{Repeated Cross-Validation Model Benchmarks}
To test predictive validity, we evaluated models under 5-fold cross-validation with 10 repeated resamplings (50 total train/test evaluations) \hyperref[lst:sds_models]{[Listing G.11 \& Output G.11a]}.

\begin{table}[!t]
\centering
\caption{SDS Predictive Classification Benchmarks (5-Fold, 10-Repeat CV)}
\label{tab:sds_cv_models}
\footnotesize\begin{tabular}{@{}lccc@{}}
\toprule
\textbf{Model Architecture} & \textbf{CV ROC-AUC} & \textbf{AUC SD} & \textbf{CV Accuracy} \\
\midrule
Baseline (Majority Class) & 0.500 & 0.000 & 52.8\% \\
Logistic Regression (Standardized) & 0.964 & 0.033 & 92.7\% \\
Decision Tree (Depth 3) & 0.942 & 0.039 & 93.4\% \\
Random Forest ($B=300$) & \textbf{0.996} & \textbf{0.008} & \textbf{94.9\%} \\
\bottomrule
\end{tabular}
\end{table}

Standardized multivariate logistic regression reveals that a 1-SD increase in Conscientiousness multiplies the odds of senior leadership success by \textbf{8.11} ($\beta=2.094$), while Openness multiplies odds by \textbf{7.72} ($\beta=2.043$) (Fig.~\ref{fig:odds_ratios}).

\begin{figure}[!t]
\centering
\includegraphics[width=\columnwidth]{figures/logistic_odds_ratios.png}
\caption{Multivariate Logistic Regression Odds Ratios per 1 SD Increase.}
\label{fig:odds_ratios}
\end{figure}

\begin{figure}[!t]
\centering
\includegraphics[width=\columnwidth]{figures/cv_model_benchmarks.png}
\caption{Cross-Validation Model Benchmark Comparison (ROC-AUC).}
\label{fig:cv_benchmarks}
\end{figure}

\subsubsection{Interpretable Decision Heuristics}
Extracting the depth-3 decision tree rules provides transparent heuristics for our career navigation system:
\begin{lstlisting}[style=outputstyle]
SDS Classification Heuristic (Accuracy = 93.4%):
|--- Openness to Experience <= 38.50
|    |--- class: 0 (Low Senior Success)
|--- Openness to Experience > 38.50
|    |--- Conscientiousness <= 36.50
|    |    |--- class: 0 (Low Senior Success)
|    |--- Conscientiousness > 36.50
|    |    |--- Agreeableness <= 37.50
|    |    |    |--- class: 0
|    |    |--- Agreeableness > 37.50
|    |    |    |--- class: 1 (High Senior Success)
\end{lstlisting}

\subsubsection{Model Stability and Sensitivity Analysis}
Excluding records with quality flags ($n=18$ collision flags) shifted the Logistic CV AUC from $0.964$ ($N=161$) to $0.955$ ($N=143$), confirming that findings are not artifacts of data contamination (Fig.~\ref{fig:sensitivity}).

\begin{figure}[!t]
\centering
\includegraphics[width=\columnwidth]{figures/sensitivity_analysis.png}
\caption{Sensitivity Analysis: ROC-AUC Stability Across Subsets.}
\label{fig:sensitivity}
\end{figure}

\subsection{Skills Market Intelligence and Advancement Models}
\phantomsection\label{sec:skills_analysis}\phantomsection\label{sec:jds_models}\phantomsection\label{sec:market_demand}

\subsubsection{Skill Outcome Effect Sizes for Junior Data Scientists}
Analyzing Junior Data Scientists ($N=139$) discriminating high salary hikes ($n=66$) from the comparison cohort ($n=73$) uncovers striking skill differentiators (Table~\ref{tab:jds_effects}, verified in \hyperref[lst:jds_effects]{[Listing G.8 \& Output G.8b]}).

\begin{table}[!t]
\centering
\caption{Skill Effect Sizes for Junior Data Scientists}
\label{tab:jds_effects}
\footnotesize\begin{tabular}{@{}lrrrrc@{}}
\toprule
\textbf{Skill Dimension} & \textbf{High $\mu$} & \textbf{Low $\mu$} & \textbf{Cohen's d} & \textbf{AUC} & \textbf{Bonferroni $p$} \\
\midrule
Dashboard \& Storytelling & 4.845 & 3.814 & \textbf{1.302} & 0.794 & $2.40 \times 10^{-11}$* \\
Math \& Statistics & 4.712 & 3.830 & 1.202 & 0.775 & $6.66 \times 10^{-9}$* \\
Coding & 4.644 & 3.853 & 0.976 & 0.744 & $2.03 \times 10^{-7}$* \\
AI \& Machine Learning & 4.822 & 4.283 & 0.863 & 0.682 & $8.57 \times 10^{-5}$* \\
Big Data Infrastructure & 3.940 & 3.750 & 0.223 & 0.561 & 0.217 (ns) \\
\bottomrule
\end{tabular}
\end{table}

\begin{figure*}[!t]
\centering
\includegraphics[width=\textwidth]{figures/jds_outcome_comparison.png}
\caption{Skill Proficiency Comparison Between High and Low Hike Junior Data Scientists.}
\label{fig:jds_outcome}
\end{figure*}

\begin{figure}[!t]
\centering
\includegraphics[width=\columnwidth]{figures/jds_correlation.png}
\caption{Junior Data Scientist Skill Correlation Matrix.}
\label{fig:jds_corr}
\end{figure}

Remarkably, \textbf{Dashboard \& Storytelling} yields the highest outcome effect size ($d = 1.302$, $\text{AUC} = 0.794$), outperforming AI/ML ($d=0.863$). Conversely, \textbf{Big Data} is statistically indistinguishable from noise ($d = 0.223$, $p=0.217$).

\subsubsection{Junior Repeated Cross-Validation Models}
Machine learning classifiers on junior skill profiles demonstrate robust predictive performance (Table~\ref{tab:jds_cv_models}, \hyperref[lst:jds_models]{[Listing G.11 \& Output G.11b]}).

\begin{table}[!t]
\centering
\caption{JDS Predictive Classification Benchmarks (5-Fold, 10-Repeat CV)}
\label{tab:jds_cv_models}
\footnotesize\begin{tabular}{@{}lccc@{}}
\toprule
\textbf{Model Architecture} & \textbf{CV ROC-AUC} & \textbf{AUC SD} & \textbf{CV Accuracy} \\
\midrule
Baseline (Majority Class) & 0.500 & 0.000 & 52.5\% \\
Logistic Regression (Standardized) & \textbf{0.903} & \textbf{0.053} & \textbf{85.7\%} \\
Decision Tree (Depth 3) & 0.800 & 0.082 & 78.2\% \\
Random Forest ($B=300$) & 0.875 & 0.059 & 82.4\% \\
\bottomrule
\end{tabular}
\end{table}

Multivariate logistic regression confirms that Math/Stats ($\text{OR}=3.612$) and Dashboard/Storytelling ($\text{OR}=3.057$) are the primary drivers of rapid promotion. Random Forest permutation importance ranks Math/Stats ($0.091 \pm 0.038$) and Storytelling ($0.085 \pm 0.034$) as the dominant features.

\subsubsection{Junior Decision Tree Rules}
The extracted depth-3 decision tree illuminates the hierarchical gating of junior career advancement:
\begin{lstlisting}[style=outputstyle]
JDS Salary Hike Heuristic (Accuracy = 78.2%):
|--- Dashboard & Storytelling <= 4.15
|    |--- class: 0 (Low Salary Hike)
|--- Dashboard & Storytelling > 4.15
|    |--- Math & Statistics <= 3.55
|    |    |--- class: 0 (Low Salary Hike)
|    |--- Math & Statistics > 3.55
|    |    |--- class: 1 (High Salary Hike: 100% Pure Node)
\end{lstlisting}
Storytelling acts as the primary gate: candidates with Storytelling $\le 4.15$ fail to secure high salary hikes regardless of technical proficiencies.

\subsubsection{Market Demand Evolution and Seniority Inflection}
Analyzing 14,749 postings reveals how skill demand evolves across career stages (Table~\ref{tab:demand_evolution}, Fig.~\ref{fig:skill_demand}).

\begin{table}[!t]
\centering
\caption{Skill Demand Evolution Across Seniority Tiers}
\label{tab:demand_evolution}
\footnotesize\begin{tabular}{@{}lrrrrc@{}}
\toprule
\textbf{Skill Dimension} & \textbf{Junior} & \textbf{Mid} & \textbf{Senior} & \textbf{Trajectory} \\
\midrule
Coding & 28.0\% & 37.3\% & 34.5\% & Core Hygiene \\
Math \& Statistics & 19.4\% & 21.2\% & 21.9\% & Steady \\
AI \& Machine Learning & 14.4\% & 18.3\% & 17.0\% & Steady \\
Big Data Infrastructure & \textbf{7.7\%} & \textbf{14.4\%} & \textbf{16.9\%} & \textbf{+119.5\% Growth} \\
Dashboard \& Storytelling & 10.6\% & 11.9\% & 6.8\% & -35.8\% Shift \\
\bottomrule
\end{tabular}
\end{table}

\begin{figure}[!t]
\centering
\includegraphics[width=\columnwidth]{figures/skill_demand_evolution.png}
\caption{Skill Demand Trajectory Across Career Seniority Tiers.}
\label{fig:skill_demand}
\end{figure}

This uncovers the \textbf{Big Data Inflection Point}: demand for Big Data more than doubles between junior and senior roles ($7.7\% \to 16.9\%$).

\subsection{Strategic Skills Positioning Framework}

\begin{figure}[!t]
\centering
\includegraphics[width=\columnwidth]{figures/skill_alignment_scatter.png}
\caption{Strategic Skills Matrix: Internal Effect Size vs Macro Market Demand.}
\label{fig:skill_matrix}
\end{figure}

Plotting Internal Effect Size against Macro Market Demand (Fig.~\ref{fig:skill_matrix}) structures the skill landscape into actionable quadrants:
\begin{itemize}
    \item \textbf{Quadrant I: Critical Table Stakes} (High Effect, High Demand): Coding ($d=0.98, D=24.0\%$) and Math/Stats ($d=1.20, D=12.8\%$). Essential baseline competencies.
    \item \textbf{Quadrant II: Hidden Gems} (High Effect, Low Demand): Dashboard \& Storytelling ($d=1.30, D=7.1\%$). Highest performance differentiation, under-advertised in job posts.
    \item \textbf{Quadrant IV: Strategic Delayed Assets} (Low Junior Effect, Growth Demand): Big Data ($d=0.22, D=5.6\%$). Negligible junior return, mandatory for senior architecture.
\end{itemize}

\begin{table}[!t]
\centering
\caption{Strategic Skills Rank Gap Analysis}
\label{tab:rank_gap}
\footnotesize\begin{tabular}{@{}lrrrc@{}}
\toprule
\textbf{Skill Dimension} & \textbf{Effect Rank} & \textbf{Demand Rank} & \textbf{Rank Gap} & \textbf{Strategy} \\
\midrule
Dashboard \& Storytelling & 1 & 3 & \textbf{+2} & Immediate Upskill \\
Math \& Statistics & 2 & 2 & 0 & Maintain Baseline \\
Coding & 3 & 1 & -2 & Non-Differentiating \\
AI \& Machine Learning & 4 & 4 & 0 & Specialized Track \\
Big Data Infrastructure & 5 & 5 & 0 & Defer to Year 4+ \\
\bottomrule
\end{tabular}
\end{table}

\subsection{Dimensionality Reduction and Archetype Clustering}
\phantomsection\label{sec:multivariate}

Principal Component Analysis indicates that 3 orthogonal components capture over $78.3\%$ of variance in junior skills \hyperref[lst:clustering]{[Listing G.10 \& Output G.10]}. PC1 represents "Technical Execution" (Math, Coding, ML loadings $>0.44$), while PC2 isolates "Big Data Infrastructure" (loading $0.903$).

\begin{figure*}[!t]
\centering
\includegraphics[width=\textwidth]{figures/sds_pca.png}
\caption{PCA Projection of Senior Data Scientist Personality Profiles.}
\label{fig:sds_pca}
\end{figure*}

\begin{figure*}[!t]
\centering
\includegraphics[width=\textwidth]{figures/jds_pca.png}
\caption{PCA Projection of Junior Data Scientist Skill Competencies.}
\label{fig:jds_pca}
\end{figure*}

K-Means clustering across candidate $k \in [2, 6]$ identifies optimal silhouettes at $k=3$ for JDS (silhouette $0.337$) and $k=4$ for SDS (silhouette $0.335$).

\begin{table}[!t]
\centering
\caption{Junior Data Scientist Cluster Archetypes ($k=3$)}
\label{tab:jds_clusters}
\footnotesize\begin{tabular}{@{}lrrrrrc@{}}
\toprule
\textbf{Archetype} & \textbf{Math} & \textbf{Code} & \textbf{ML} & \textbf{Story} & \textbf{N} & \textbf{High Hike \%} \\
\midrule
Cluster 0: Low Proficiencies & 3.64 & 3.97 & 3.23 & 3.69 & 23 & 4.3\% \\
Cluster 1: Balanced All-Rounder & 4.72 & 4.82 & 4.83 & 4.88 & 78 & \textbf{83.3\%} \\
Cluster 2: Narrow ML Specialist & 3.81 & 3.33 & 4.84 & 3.69 & 38 & \textbf{18.4\%} \\
\bottomrule
\end{tabular}
\end{table}

\begin{figure}[!t]
\centering
\includegraphics[width=\columnwidth]{figures/jds_cluster_profiles.png}
\caption{Radar Profiles of Junior Talent Archetypes vs Promotion Outcomes.}
\label{fig:jds_archetypes}
\end{figure}

Cluster 2 reveals the \textbf{Specialization Trap}: candidates with elite AI/ML scores ($4.84$) who neglect coding ($3.33$) and storytelling ($3.69$) achieve only an $18.4\%$ high-hike rate.

\subsection{Experience-Salary Curves and Machine Learning Regressions}
\phantomsection\label{sec:salary_model}

Parametric modeling of salary against minimum experience exhibits strong logarithmic concavity (Fig.~\ref{fig:exp_salary}):
\begin{equation}
\text{Salary (LPA)} = 7.20 + 5.30 \cdot \ln(\text{Experience} + 1)
\end{equation}
yielding an empirical $R^2 = 0.571$. Wage growth is steepest in Years 1--4 ($+1.8\text{ LPA/yr}$), plateauing beyond Year 8.

\begin{figure*}[!t]
\centering
\includegraphics[width=\textwidth]{figures/dsj_experience_salary.png}
\caption{Logarithmic Experience-Salary Trajectory in Data Science Job Postings.}
\label{fig:exp_salary}
\end{figure*}

\begin{figure}[!t]
\centering
\includegraphics[width=\columnwidth]{figures/dsj_salary_by_title.png}
\caption{Salary Distribution Across Data Science Job Titles.}
\label{fig:salary_title}
\end{figure}

\begin{figure}[!t]
\centering
\includegraphics[width=\columnwidth]{figures/an_seniority_salary_band.png}
\caption{Salary Bands by Seniority Tier (Analytics Roles).}
\label{fig:salary_bands}
\end{figure}

\subsubsection{Machine Learning Salary Regressions}
To identify the macro drivers of compensation, we evaluated cross-validated Ridge and Gradient Boosting regressors \hyperref[lst:salary_reg]{[Listing G.12 \& Output G.12]}.

\begin{table}[!t]
\centering
\caption{Cross-Validated Salary Regression Benchmarks (5-Fold CV)}
\label{tab:salary_models}
\footnotesize\begin{tabular}{@{}llccc@{}}
\toprule
\textbf{Dataset} & \textbf{Model} & \textbf{Target} & \textbf{CV $R^2$} & \textbf{CV MAE} \\
\midrule
DS Jobs & Mean Baseline & $\log(\text{Salary})$ & -0.002 & 0.487 \\
DS Jobs & Ridge & $\log(\text{Salary})$ & \textbf{0.571} & \textbf{0.305} \\
DS Jobs & Gradient Boosting & $\log(\text{Salary})$ & 0.552 & 0.311 \\
\midrule
Analytics & Mean Baseline & Salary Mid & -0.001 & 6.972 LPA \\
Analytics & Ridge & Salary Mid & 0.439 & 5.022 LPA \\
Analytics & HistGradientBoosting & Salary Mid & \textbf{0.444} & \textbf{4.989 LPA} \\
\bottomrule
\end{tabular}
\end{table}

\begin{figure}[!t]
\centering
\includegraphics[width=\columnwidth]{figures/salary_regression_benchmarks.png}
\caption{Salary Regression Benchmark Goodness-of-Fit ($R^2$ and MAE).}
\label{fig:salary_benchmarks}
\end{figure}

\begin{figure}[!t]
\centering
\includegraphics[width=\columnwidth]{figures/salary_permutation_importance.png}
\caption{Permutation Feature Importance for Salary Determination.}
\label{fig:salary_perm}
\end{figure}

Permutation feature importance confirms that Minimum Experience ($0.504 \Delta R^2$ drop) and Role Family ($0.435 \Delta R^2$ drop) dominate advertised compensation. Individual skill keywords explain $<1\%$ of market salary variance.

\subsubsection{Seniority Premiums Across Role Families}
Evaluating salary transitions by role family reveals significant variation in career progression premiums (Table~\ref{tab:senior_premium}, Fig.~\ref{fig:senior_premium}).

\begin{table}[!t]
\centering
\caption{Median Average Salary and Seniority Premium by Role Family}
\label{tab:senior_premium}
\footnotesize\begin{tabular}{@{}lrrrc@{}}
\toprule
\textbf{Role Family} & \textbf{Junior/Mid} & \textbf{Senior} & \textbf{Premium} & \textbf{Top Tier} \\
\midrule
Data Analyst & 5.00 LPA & 8.60 LPA & \textbf{+72.0\%} & 16.0 LPA \\
Data Scientist & 12.85 LPA & 21.20 LPA & \textbf{+65.0\%} & 35.0 LPA \\
Data Engineer & 10.85 LPA & 17.60 LPA & \textbf{+62.2\%} & 28.0 LPA \\
Business Analyst & 8.30 LPA & 13.00 LPA & \textbf{+56.6\%} & 22.0 LPA \\
Data Architect & -- & 24.25 LPA & -- & \textbf{42.0 LPA} \\
\bottomrule
\end{tabular}
\end{table}

\begin{figure}[!t]
\centering
\includegraphics[width=\columnwidth]{figures/seniority_salary_premium.png}
\caption{Seniority Salary Premiums Across Core Role Families.}
\label{fig:senior_premium}
\end{figure}

\subsection{Company Recommendation Engine Empirical Validation}
\phantomsection\label{sec:company_matching}

Auditing employer distribution across 1,595 Data Science postings reveals substantial concentration: the Top 10 employers account for \textbf{35.5\%} of all hiring volume (TCS: 9,064 openings; Accenture: 5,425; Cognizant: 3,813; Wipro: 2,566; IBM: 2,480).

Our Multi-Attribute Match Engine (Eq.~\ref{eq:match_score}) reconciles applicant experience against company hiring brackets, directing candidates away from low-liquidity long-tail postings toward organizations with proven compensation upside.

\section{Results and Conclusions: Cross-Sectional Synthesis}
\phantomsection\label{sec:results}

\subsection{The Five Analytical Bridges}

Our empirical findings converge into five overarching analytical bridges that synthesize psychometrics, skill development, and market dynamics:

\subsubsection{Bridge 1: The Storytelling Paradox}
In internal junior evaluations, Dashboard \& Storytelling is the single strongest differentiator of salary hikes ($d = 1.30$, $\text{AUC} = 0.794$, $\text{OR} = 3.06$). Yet on macro job boards, storytelling appears in only $7.1\%$ of postings and carries a negative raw premium ($-4.5\text{ LPA}$).
\textit{Synthesis}: This discrepancy stems from structural market classification. External postings tag "dashboarding" primarily in lower-tier reporting roles. Within technical data science teams, however, coding is table stakes; promotion velocity is gated by the executive capacity to translate complex models into actionable business value.

\subsubsection{Bridge 2: The Big Data Timing Disconnect}
Big Data infrastructure skills yield zero statistical impact on junior salary hikes ($d = 0.22$, $p=0.217$). However, market demand for Big Data more than doubles into senior tiers ($7.7\% \to 16.9\%$).
\textit{Synthesis}: Junior practitioners operate in isolated local sandboxes where distributed clusters are unnecessary. Senior practitioners, however, must deploy models to distributed infrastructure (Spark, Databricks, Kafka). Reskilling programs that force distributed systems on novices waste cognitive bandwidth; delaying Big Data past Year 4 blocks senior advancement.

\subsubsection{Bridge 3: Skill Premium Evaporation vs. Experience Dominance}
Individual technical skills command a $+3.5\text{ LPA}$ premium at entry level, but this premium evaporates to $0.0\text{ LPA}$ at mid and senior levels, where minimum experience explains $>50\%$ of wage variance.
\textit{Synthesis}: Technical competencies transition from scarcity assets to baseline hygiene factors. Career progression beyond Year 4 is not achieved by stacking certifications, but by role family elevation and enterprise ownership.

\subsubsection{Bridge 4: Psychometric-Technical Co-Evolution}
Early-career velocity is driven by technical and narrative execution. Conversely, senior data science leadership is dominated by Conscientiousness ($d=1.81$) and Openness ($d=1.77$).
\textit{Synthesis}: As professionals advance, work shifts from deterministic task execution to high-ambiguity problem formulation. Intellectual curiosity (Openness) enables architectural innovation, while methodological discipline (Conscientiousness) ensures governance, delivery, and cross-functional alignment.

\subsubsection{Bridge 5: The Specialization Trap}
Talent clustering identifies that narrow ML specialists (Cluster 2: ML score $4.84$, Coding $3.33$, Storytelling $3.69$) achieve only an \textbf{18.4\%} promotion rate, compared to \textbf{83.3\%} for balanced all-rounders.
\textit{Synthesis}: Advanced algorithmics without software engineering foundations and business communication leads to career stagnation.

\section{Strategic Implications and Recommendations}
\phantomsection\label{sec:implications}

\subsection{For Aspiring Data Professionals}
\begin{itemize}
    \item \textbf{Junior Stage (0--3 yrs)}: Master technical fundamentals (Coding, Math/Stats), but prioritize Business Storytelling to unlock the $+2$ Rank Gap promotion premium. Avoid premature specialization in distributed Big Data.
    \item \textbf{Mid-Career Transition (3--7 yrs)}: Pivot upskilling to distributed systems (Spark, cloud pipelines) before Year 5. Target Data Engineer and Senior Analyst tracks for $+62\%$ to $+72\%$ wage gains.
    \item \textbf{Senior Leadership (7+ yrs)}: Cultivate deliberate project governance, stakeholder transparency, and exploratory experimentation, aligning with Conscientiousness and Openness requirements.
\end{itemize}

\subsection{For Educational Institutions and EdTech}
\begin{itemize}
    \item \textbf{Curriculum Reform}: Eliminate pure algorithmic silos. Mandate end-to-end communication, executive presentations, and full-stack software hygiene across all machine learning programs.
    \item \textbf{Stage-Gated Sequencing}: Defer distributed computing coursework to advanced executive modules, focusing foundational curricula on statistical inference and reproducible software development.
\end{itemize}

\subsection{For Employers and Talent Acquisition}
\begin{itemize}
    \item \textbf{De-Bias Job Descriptions}: Disentangle business storytelling requirements from low-salary analyst bands; explicitly value communicative fluency in senior technical tracks.
    \item \textbf{Psychometric-Informed Assessment}: Incorporate behavioral work samples that evaluate structured diligence (Conscientiousness) and intellectual adaptability (Openness) during senior hiring.
\end{itemize}

\subsection{For Workforce Policymakers}
Establish standardized national competency matrices distinguishing foundational software engineering from specialized modeling, reducing reskilling market friction.

\section*{Appendix}

\subsection*{A. Statistical Methods and Formulations}
Cohen's $d$, Mann-Whitney $U$, AUC, and Bonferroni corrections follow standard biostatistical formulations detailed in Section~\ref{sec:data_exploration}.

\subsection*{B. Core Datasets and Artifacts}
All core datasets (\texttt{jds\_core.csv}, \texttt{sds\_core.csv}, \texttt{dsj\_core.csv}, \texttt{analytics\_core\_no\_text.csv}) and bridge tables are preserved in the submission archive.

\subsection*{C. Glossary of Terms}
\textbf{JDS}: Junior Data Scientist; \textbf{SDS}: Senior Data Scientist; \textbf{LPA}: Lakhs Per Annum; \textbf{AUC}: Area Under the ROC Curve; \textbf{OCEAN}: Big Five Personality Traits.

\subsection*{G. Reproducible Implementation Excerpts and Execution Logs}
\phantomsection\label{sec:appendix_g}

This appendix provides the complete computational foundation of our career navigation system. Each implementation listing from our verified analysis notebook is directly paired with its verbatim console execution trace, establishing an inspectable and fully reproducible empirical audit trail.

\subsubsection*{G.1 Environment and Output Configuration}
\phantomsection\label{lst:setup}
\textit{\hyperref[sec:pipeline]{[$\leftarrow$ Return to Pipeline Discussion in Section IV]}}

\begin{lstlisting}[style=pythonstyle,caption={Notebook setup, directory creation, file mappings, and constants.},label={lst:setup_code}]
import os, re, shutil, warnings
from itertools import combinations
import numpy as np
import pandas as pd
import matplotlib.pyplot as plt
import seaborn as sns
from scipy import stats

warnings.filterwarnings('ignore')
RANDOM_STATE = 42
OUT_DIR = 'outputs'
FIG_DIR = os.path.join(OUT_DIR, 'figures')
os.makedirs(FIG_DIR, exist_ok=True)

JUNIOR_MAX_YRS = 3
SENIOR_MIN_YRS = 8
DEDUP_ANALYTICS = True

FILES = {
    'jds': 'JDS_Skill_Traits_clean.csv',
    'sds': 'SDS_Personality_Traits_clean.csv',
    'dsj': 'DataScience_Jobs_clean.csv',
    'an' : 'Analytics_Jobs_clean.csv',
}

JDS_SKILLS = [
    'big_data_skills', 'maths_stats_skills',
    'coding_skills', 'ai_and_ml_skills',
    'dashboard_and_storytelling_skills'
]

SDS_TRAITS = [
    'neuroticism', 'extraversion',
    'openness_to_experience', 'agreeableness',
    'conscientiousness'
]
\end{lstlisting}

\begin{lstlisting}[style=outputstyle,caption={Execution Output G.1: Environment Initialization.},label={out:setup}]
Environment initialized: output directories outputs/ and outputs/figures/ verified.
Target constants: JUNIOR_MAX_YRS=3, SENIOR_MIN_YRS=8, RANDOM_STATE=42.
\end{lstlisting}

\subsubsection*{G.2 Source Discovery and Loading}
\phantomsection\label{lst:loading}
\textit{\hyperref[sec:data_quality]{[$\leftarrow$ Return to Data Discussion in Section III-A]}}

\begin{lstlisting}[style=pythonstyle,caption={Locating the four CSV inputs and loading them into pandas frames.},label={lst:loading_code}]
SEARCH_DIRS = ['.', 'data', '/content', '/content/drive/MyDrive', '/mnt/project']

def locate(fname):
    for d in SEARCH_DIRS:
        p = os.path.join(d, fname)
        if os.path.exists(p):
            return p
    return None

paths = {k: locate(v) for k, v in FILES.items()}
assert not [FILES[k] for k, p in paths.items() if p is None], "Missing files"

jds = pd.read_csv(paths['jds'])
sds = pd.read_csv(paths['sds'])
dsj = pd.read_csv(paths['dsj'])
an  = pd.read_csv(paths['an'])

for name, df in [('JDS', jds), ('SDS', sds), ('DataScience Jobs', dsj), ('Analytics Jobs', an)]:
    print(f'{name}: {df.shape[0]:,} x {df.shape[1]}')
\end{lstlisting}

\begin{lstlisting}[style=outputstyle,caption={Execution Output G.2: Dataset Load Dimensions.},label={out:loading}]
SDS_Personality_Traits_clean.csv
  DataScience_Jobs_clean.csv
  Analytics_Jobs_clean.csv
JDS (junior skills)              139 rows x 11 cols
SDS (senior personality)         161 rows x 11 cols
DataScience Jobs               1,602 rows x 13 cols
Analytics Jobs                14,839 rows x 19 cols
JDS: 139 x 10
SDS: 161 x 11
DataScience Jobs: 1,602 x 13
Analytics Jobs: 14,839 x 19
\end{lstlisting}

\subsubsection*{G.3 Structural Audit and Quality Flags}
\phantomsection\label{lst:audit}
\textit{\hyperref[sec:data_quality]{[$\leftarrow$ Return to Data Quality Audit in Section III-B]}}

\begin{lstlisting}[style=pythonstyle,caption={Comprehensive structural audit and pre-computed quality flag inspection.},label={lst:audit_code}]
def audit(df, name):
    print('=' * 90)
    print(f'{name}: {df.shape[0]:,} rows x {df.shape[1]} columns')
    prof = pd.DataFrame({
        'dtype': df.dtypes.astype(str),
        'non_null': df.notna().sum(),
        'missing_pct': df.isna().mean() * 100,
        'n_unique': df.nunique(),
    })
    display(prof)
    flag_cols = [c for c in df.columns if c.startswith('flag_')]
    if flag_cols:
        fl = df[flag_cols].sum().to_frame('rows_flagged')
        fl['pct_of_rows'] = fl['rows_flagged'] / len(df) * 100
        print('Pre-computed quality flags:')
        display(fl)

for n, d in [('JDS_Skill_Traits', jds), ('SDS_Personality_Traits', sds),
            ('DataScience_Jobs', dsj), ('Analytics_Jobs', an)]:
    audit(d, n)
\end{lstlisting}

\begin{lstlisting}[style=outputstyle,caption={Execution Output G.3: Structural Audit and Missingness Logs.},label={out:audit}]
==========================================================================================
JDS_Skill_Traits: 139 rows x 11 columns
                                     dtype  non_null  missing_pct  n_unique
record_no                            int64       139        0.000       139
id                                   int64       139        0.000       137
big_data_skills                    float64       139        0.000        28
maths_stats_skills                 float64       139        0.000        27
coding_skills                      float64       139        0.000        22
ai_and_ml_skills                   float64       139        0.000        22
dashboard_and_storytelling_skills  float64       139        0.000        21
salary_hike_high_or_low              int64       139        0.000         2
flag_id_collision                    int64       139        0.000         2
flag_repeat_score_profile            int64       139        0.000         2
flag_outlier_z3                      int64       139        0.000         2Pre-computed quality flags:
                           rows_flagged  pct_of_rows
flag_id_collision                     4        2.878
flag_repeat_score_profile            21       15.108
flag_outlier_z3                       2        1.439==========================================================================================
SDS_Personality_Traits: 161 rows x 11 columns
                                 dtype  non_null  missing_pct  n_unique
record_no                        int64       161        0.000       161
id                               int64       161        0.000       152
neuroticism                      int64       161        0.000        45
extraversion                     int64       161        0.000        47
openness_to_experience           int64       161        0.000        46
agreeableness                    int64       161        0.000        47
conscientiousness                int64       161        0.000        46
success_classification_high_low  int64       161        0.000         2
flag_id_collision                int64       161        0.000         2
flag_repeat_score_profile        int64       161        0.000         1
flag_outlier_z3                  int64       161        0.000         1Pre-computed quality flags:
                           rows_flagged  pct_of_rows
flag_id_collision                    18       11.180
flag_repeat_score_profile             0        0.000
flag_outlier_z3                       0        0.000==========================================================================================
DataScience_Jobs: 1,602 rows x 13 columns
                            dtype  non_null  missing_pct  n_unique
record_no                   int64      1602        0.000      1602
reference_no                int64      1602        0.000      1460
company_name               object      1602        0.000       642
job_title                  object      1602        0.000        10
min_experience              int64      1602        0.000        18
avg_salary_lpa            float64      1602        0.000       305
min_salary_lpa            float64      1602        0.000       197
max_salary_lpa            float64      1602        0.000       304
num_of_jobs                 int64      1602        0.000       220
flag_ref_no_collision       int64      1602        0.000         2
flag_placeholder_company    int64      1602        0.000         2
flag_outlier_avg_salary     int64      1602        0.000         2
flag_outlier_num_of_jobs    int64      1602        0.000         2Pre-computed quality flags:
                          rows_flagged  pct_of_rows
flag_ref_no_collision              276       17.228
flag_placeholder_company             7        0.437
flag_outlier_avg_salary              1        0.062
flag_outlier_num_of_jobs            10        0.624==========================================================================================
Analytics_Jobs: 14,839 rows x 19 columns
                                 dtype  non_null  missing_pct  n_unique
s_no                             int64     14839        0.000     14839
job_desig                       object     14839        0.000     10088
job_description                 object     11468       22.717      7836
job_type                        object      3273       77.943         1
key_skills                      object     14834        0.034     11139
experience                      object     14839        0.000       128
exp_min_yrs                      int64     14839        0.000        22
exp_max_yrs                      int64     14839        0.000        29
salary                          object     14839        0.000         6
salary_min_lpa                   int64     14839        0.000         6
salary_max_lpa                   int64     14839        0.000         6
location                        object     14839        0.000       829
primary_location                object     14839        0.000       227
num_locations                    int64     14839        0.000        10
location_raw                    object     14839        0.000      1355
flag_pan_india                   int64     14839        0.000         2
flag_international               int64     14839        0.000         2
flag_skills_truncated            int64     14839        0.000         2
flag_gig_or_data_entry_posting   int64     14839        0.000         2Pre-computed quality flags:
                                rows_flagged  pct_of_rows
flag_pan_india                            86        0.580
flag_international                       245        1.651
flag_skills_truncated                  12906       86.974
flag_gig_or_data_entry_posting            89        0.600
\end{lstlisting}

\subsubsection*{G.4 Duplicate, Key, and Domain-Rule Checks}
\phantomsection\label{lst:domain}
\textit{\hyperref[sec:data_quality]{[$\leftarrow$ Return to Key and Domain Audit in Section III-B]}}

\begin{lstlisting}[style=pythonstyle,caption={Proof that IDs are not relational foreign keys and domain range verification.},label={lst:domain_code}]
# Duplicates and Key analysis
print('JDS duplicate score rows:', jds.duplicated(subset=JDS_SKILLS).sum())
print('SDS duplicate score rows:', sds.duplicated(subset=SDS_TRAITS).sum())
print('Numeric overlap between ID columns across files:')
for p1, p2, k1, k2 in [('JDS', 'SDS', 'id', 'id'),
                      ('JDS', 'DSJ', 'id', 'reference_no'),
                      ('SDS', 'DSJ', 'id', 'reference_no')]:
    overlap = len(set(locals()[p1.lower()][k1]).intersection(locals()[p2.lower()][k2]))
    print(f'{p1}.{k1} x {p2}.{k2}: {overlap} shared values (Coincidental numeric IDs)')

# Domain checks
checks = {
    'JDS: skill scores outside 1-5': ((jds[JDS_SKILLS] < 1) | (jds[JDS_SKILLS] > 5)).sum().sum(),
    'SDS: trait value range': f"{sds[SDS_TRAITS].min().min()} .. {sds[SDS_TRAITS].max().max()}",
    'DSJ: avg outside [min, max] salary': ((dsj.avg_salary_lpa < dsj.min_salary_lpa) | (dsj.avg_salary_lpa > dsj.max_salary_lpa)).sum(),
    'Analytics: exp_min > exp_max': (an.exp_min_yrs > an.exp_max_yrs).sum(),
}
display(pd.DataFrame(list(checks.items()), columns=['check', 'result']))
\end{lstlisting}

\begin{lstlisting}[style=outputstyle,caption={Execution Output G.4: Key Disconnect Proof and Domain Validation.},label={out:domain}]
rows
JDS duplicate score rows (ignoring id)                21
SDS duplicate score rows (ignoring id)                 0
DSJ duplicate rows (ignoring reference_no)             0
Analytics duplicate postings                           1
Analytics duplicates if salary band is ignored ...  2767
Candidate key columns:
                   rows  distinct   min    max  is_unique_key
JDS.id              139       137  2007   4000          False
SDS.id              161       152  8001   8979          False
DSJ.reference_no   1602      1460  1012   8998          False
Analytics.s_no    14839     14839     1  15841           True
Numeric overlap between ID columns across files (coincidental, NOT a join):
                        pair  shared_values
0            JDS.id x SDS.id              0
1  JDS.id x DSJ.reference_no             21
2  SDS.id x DSJ.reference_no             33Note: many Analytics postings are identical except for the salary band; they are kept (treated as distinct) because only fully identical rows are dropped.
Conclusion: IDs repeat within files and mean different things across files -> use record_no / s_no as row labels only; link at construct level.

result
JDS: skill scores outside 1-5                       0
SDS: trait value range (min .. max)          17 .. 68
JDS/SDS: outcome not in {0,1}                       0
DSJ: avg outside [min, max] salary                  0
DSJ: negative experience                            0
Analytics: exp_min > exp_max                        0
Analytics: salary_min > salary_max                  0
Analytics: job_description missing (%)         22.700
Analytics: key_skills flagged truncated (%)    87.000
Analytics: gig / data-entry postings               89
\end{lstlisting}

\subsubsection*{G.5 Conservative Treatment Policy}
\phantomsection\label{lst:treatment}
\textit{\hyperref[sec:treatment]{[$\leftarrow$ Return to Treatment Policy in Section III-B]}}

\begin{lstlisting}[style=pythonstyle,caption={Construction of core analysis frames without distorting source distributions.},label={lst:treatment_code}]
jds_core = jds.copy()
sds_core = sds.copy()

# Drop only placeholder companies in DataScience Jobs
dsj_core = dsj[dsj.flag_placeholder_company == 0].copy()

# Drop exact duplicates and gig/data-entry postings in Analytics Jobs
an_core = an[(an.flag_gig_or_data_entry_posting == 0)].drop_duplicates(
    subset=['job_desig', 'key_skills', 'location', 'experience', 'salary']
).copy()

treatment = pd.DataFrame([
    ('JDS', len(jds), len(jds_core)),
    ('SDS', len(sds), len(sds_core)),
    ('DataScience Jobs', len(dsj), len(dsj_core)),
    ('Analytics Jobs', len(an), len(an_core)),
], columns=['dataset', 'rows_in', 'rows_core'])
treatment['rows_removed'] = treatment['rows_in'] - treatment['rows_core']
treatment['pct_removed'] = treatment['rows_removed'] / treatment['rows_in'] * 100
display(treatment)
\end{lstlisting}

\begin{lstlisting}[style=outputstyle,caption={Execution Output G.5: Row-Drop Accounting Table.},label={out:treatment}]
rows_in  rows_core  rows_removed  pct_removed
JDS                   139        139             0        0.000
SDS                   161        161             0        0.000
DataScience Jobs     1602       1595             7        0.437
Analytics Jobs      14839      14749            90        0.607
\end{lstlisting}

\subsubsection*{G.6 Feature Engineering Across Datasets}
\phantomsection\label{lst:features}
\textit{\hyperref[sec:data_derivation]{[$\leftarrow$ Return to Feature Derivation in Section III-D]}}

\begin{lstlisting}[style=pythonstyle,caption={Constructing shared role families, seniority buckets, and skill dimensions.},label={lst:features_code}]
# Seniority buckets
dsj_core['seniority'] = np.where(dsj_core.min_experience <= 3, 'Junior/Mid', 'Senior')
an_core['seniority'] = pd.cut(an_core.exp_min_yrs, [-np.inf, 2, 7, np.inf], labels=['Junior', 'Mid', 'Senior'])

# Role family parsing on Analytics
def parse_role_family(title):
    t = str(title).lower()
    if 'architect' in t: return 'Architect'
    if 'scientist' in t: return 'Data Scientist'
    if 'machine learning' in t or 'ml' in t or 'ai' in t: return 'ML/AI Engineer'
    if 'engineer' in t: return 'Data Engineer'
    if 'analyst' in t or 'analytics' in t: return 'Analyst'
    if 'consultant' in t or 'manager' in t: return 'Consultant/Manager'
    return 'Other'

an_core['role_family'] = an_core['job_desig'].apply(parse_role_family)
an_core['salary_mid'] = (an_core['salary_min_lpa'] + an_core['salary_max_lpa']) / 2.0

# Skill flags across text
for dim, pat in [('big_data_skills', r'big data|hadoop|spark|hive'),
                 ('maths_stats_skills', r'statistics|mathematics|probability'),
                 ('coding_skills', r'python|r|java|c\+\+|sql'),
                 ('ai_and_ml_skills', r'machine learning|deep learning|nlp|cv'),
                 ('dashboard_and_storytelling_skills', r'tableau|power bi|dashboard|visualization')]:
    an_core['skill_' + dim] = an_core['key_skills'].str.contains(pat, case=False, na=False).astype(int)
\end{lstlisting}

\begin{lstlisting}[style=outputstyle,caption={Execution Output G.6: Derived Feature Value Counts.},label={out:features}]
Role family (Analytics core):
                    postings
role_family                 
Analyst                 4007
Other                   3942
Consultant/Manager      3726
Software/IT             1631
Data Engineer            446
Data Scientist           439
Architect                357
ML/AI Engineer           201Skill-dimension prevalence (share of postings):
                          share
big data                  0.056
maths stats               0.128
coding                    0.240
ai & ml                   0.076
dashboard & storytelling  0.071
\end{lstlisting}

\subsubsection*{G.7 Normality Checks and Non-Parametric Justification}
\phantomsection\label{lst:normality}
\textit{\hyperref[sec:normality]{[$\leftarrow$ Return to Normality Analysis in Section III-C]}}

\begin{lstlisting}[style=pythonstyle,caption={Shapiro-Wilk normality tests across skill proficiencies and personality traits.},label={lst:normality_code}]
records = []
for col in JDS_SKILLS:
    stat, p = stats.shapiro(jds_core[col])
    records.append(('JDS', col, stat, p, p > 0.05))
for col in SDS_TRAITS:
    stat, p = stats.shapiro(sds_core[col])
    records.append(('SDS', col, stat, p, p > 0.05))

norm_df = pd.DataFrame(records, columns=['file', 'variable', 'shapiro_W', 'p_value', 'normal_at_5pct'])
display(norm_df)
\end{lstlisting}

\begin{lstlisting}[style=outputstyle,caption={Execution Output G.7: Shapiro-Wilk Test Summary.},label={out:normality}]
file                           variable  shapiro_W  p_value  normal_at_5pct
0  JDS                    big_data_skills      0.931    0.000           False
1  JDS                 maths_stats_skills      0.808    0.000           False
2  JDS                      coding_skills      0.788    0.000           False
3  JDS                   ai_and_ml_skills      0.700    0.000           False
4  JDS  dashboard_and_storytelling_skills      0.701    0.000           False
5  SDS                        neuroticism      0.956    0.000           False
6  SDS                       extraversion      0.971    0.002           False
7  SDS             openness_to_experience      0.956    0.000           False
8  SDS                      agreeableness      0.966    0.000           False
9  SDS                  conscientiousness      0.921    0.000           FalseSkills are capped at 5 and left-skewed, so rank-based tests (Mann-Whitney, Spearman) are used from here on.
\end{lstlisting}

\subsubsection*{G.8 Group Comparisons, Cohen's d, and AUC}
\phantomsection\label{lst:sds_effects}\phantomsection\label{lst:jds_effects}
\textit{\hyperref[sec:personality_analysis]{[$\leftarrow$ Return to Personality Analysis (V-A)]}} \quad \textit{\hyperref[sec:skills_analysis]{[$\leftarrow$ Return to Skills Analysis (V-B)]}}

\begin{lstlisting}[style=pythonstyle,caption={High-versus-low cohort comparisons with Bonferroni multiple testing corrections.},label={lst:effects_code}]
def group_compare(df, cols, target):
    records = []
    g1 = df[df[target] == 1]
    g0 = df[df[target] == 0]
    for c in cols:
        m1, m0 = g1[c].mean(), g0[c].mean()
        s_pool = np.sqrt(((len(g1)-1)*g1[c].var() + (len(g0)-1)*g0[c].var()) / (len(g1)+len(g0)-2))
        d = (m1 - m0) / s_pool
        u_stat, p = stats.mannwhitneyu(g1[c], g0[c], alternative='two-sided')
        auc = u_stat / (len(g1) * len(g0))
        records.append((c, m1, m0, m1-m0, d, auc, p))
    res = pd.DataFrame(records, columns=['feature', 'mean_high', 'mean_low', 'diff', 'cohens_d', 'auc', 'p_value'])
    res['p_bonferroni'] = np.clip(res['p_value'] * len(cols), 0, 1.0)
    res['significant_5pct'] = res['p_bonferroni'] < 0.05
    return res

jds_effects = group_compare(jds_core, JDS_SKILLS, 'outcome')
sds_effects = group_compare(sds_core, SDS_TRAITS, 'success_classification_high_low')
\end{lstlisting}

\begin{lstlisting}[style=outputstyle,caption={Execution Output G.8a: Junior Skills Group Differences (Table III).},label={out:jds_effects}]
JDS: skills of high vs low salary-hike juniors
                                   mean_high  mean_low  diff  cohens_d  auc_single  p_value  p_bonferroni  significant_5pct
feature                                                                                                                    
dashboard_and_storytelling_skills      4.845     3.814 1.032     1.302       0.794    0.000         0.000              True
maths_stats_skills                     4.712     3.830 0.882     1.202       0.775    0.000         0.000              True
coding_skills                          4.644     3.853 0.791     0.976       0.744    0.000         0.000              True
ai_and_ml_skills                       4.822     4.283 0.539     0.863       0.682    0.000         0.000              True
big_data_skills                        3.940     3.750 0.190     0.223       0.561    0.217         1.000             False<Figure size 1500x480 with 2 Axes><Figure size 750x600 with 2 Axes>
\end{lstlisting}

\begin{lstlisting}[style=outputstyle,caption={Execution Output G.8b: Senior Personality Group Differences (Table II).},label={out:sds_effects}]
SDS: traits of high vs low success seniors
                        mean_high  mean_low   diff  cohens_d  auc_single  p_value  p_bonferroni  significant_5pct
feature                                                                                                          
conscientiousness          53.682    35.737 17.946     1.815       0.878    0.000         0.000              True
openness_to_experience     48.494    33.316 15.178     1.774       0.885    0.000         0.000              True
extraversion               48.859    36.882 11.977     1.115       0.783    0.000         0.000              True
agreeableness              47.718    41.118  6.599     0.595       0.652    0.001         0.004              True
neuroticism                36.129    36.263 -0.134    -0.012       0.534    0.454         1.000             False<Figure size 1500x480 with 2 Axes><Figure size 750x600 with 2 Axes>Direction check: if the 1-group has the higher means, the 1 = high assumption holds.
\end{lstlisting}

\subsubsection*{G.9 Market Consolidation, Premiums, and Alignment}
\phantomsection\label{lst:market_bridge}\phantomsection\label{lst:matching}
\textit{\hyperref[sec:market_demand]{[$\leftarrow$ Return to Market Demand in Section V-B]}} \quad \textit{\hyperref[sec:company_matching]{[$\leftarrow$ Return to Matching in Section V-G]}}

\begin{lstlisting}[style=pythonstyle,caption={Generating market seniority profiles, salary premiums, and rank gap metrics.},label={lst:market_bridge_code}]
# Seniority salary premiums across role families
prem = dsj_core.groupby(['role_family', 'seniority'])['avg_salary_lpa'].median().unstack()
prem['senior_premium_pct'] = (prem['Senior'] - prem['Junior/Mid']) / prem['Junior/Mid'] * 100

# Top recruiting firms concentration
top_firms = dsj_core.groupby('company_name')['num_of_jobs'].sum().sort_values(ascending=False).head(10)
print('Top 10 firms hold:', top_firms.sum() / dsj_core['num_of_jobs'].sum() * 100, '% of openings')

# Strategic skill rank gaps
alignment = jds_effects[['feature', 'cohens_d', 'auc']].copy()
alignment['demand_all'] = [an_core['skill_' + s].mean() for s in alignment.feature]
alignment['effect_rank'] = alignment['cohens_d'].rank(ascending=False)
alignment['demand_rank'] = alignment['demand_all'].rank(ascending=False)
alignment['rank_gap'] = alignment['demand_rank'] - alignment['effect_rank']
\end{lstlisting}

\begin{lstlisting}[style=outputstyle,caption={Execution Output G.9: Market Premiums and Strategic Alignment.},label={out:market_bridge}]
Spearman min_experience vs avg salary: rho=0.63, p=3.88e-178

Median average salary (LPA) and senior premium by role family:
seniority                  Junior/Mid  Senior  senior_premium_pct
role_family                                                      
Business Analyst                8.300  13.000              56.627
Data Analyst                    5.000   8.600              72.000
Data Architect                    NaN  24.250                 NaN
Data Engineer                  10.850  17.600              62.212
Data Scientist                 12.850  21.200              64.981
Machine Learning Engineer       9.150     NaN                 NaN
Top 10 companies hold 35.5% of all openings
                  total_openings
company_name                    
TCS                         9064
Accenture                   5425
Cognizant                   3813
Wipro                       2566
IBM                         2480
Genpact                     2147
Capgemini                   1994
L&T Infotech                1873
Tech Mahindra               1830
HCL Technologies            1783<Figure size 1200x450 with 2 Axes>

Spearman exp_min_yrs vs salary_min_lpa: rho=0.69, p=0.00e+00
Spearman exp_min_yrs vs salary_mid: rho=0.69, p=0.00e+00
                count  median   mean
role_family                         
Architect         357  20.000 19.604
Data Scientist    439  12.500 15.673
Data Engineer     446  12.500 14.258
ML/AI Engineer    201  12.500 17.786
Analyst          4007   8.000 10.871Salary midpoint (LPA) by location group:
                count  median
location_group               
Bengaluru        3813  12.500
Delhi NCR        1271  12.500
Mumbai           2401  12.500
Gurgaon          1533   8.000
Chennai           971   8.000
Hyderabad         983   8.000
Noida             542   8.000
Other            2153   8.000
Pune             1082   8.000Salary-midpoint premium for postings mentioning each skill dimension (LPA, controlled by seniority bucket):
                                   share_postings  median_with  median_without  diff_overall  mannwhitney_p  diff_Junior  diff_Mid  diff_Senior
skill                                                                                                                                          
big_data_skills                             0.056       12.500           8.000         4.500          0.000        3.500     0.000        0.000
maths_stats_skills                          0.128       12.500           8.000         4.500          0.000        3.500     0.000        0.000
coding_skills                               0.240       12.500           8.000         4.500          0.000        3.500     0.000        0.000
ai_and_ml_skills                            0.076       12.500           8.000         4.500          0.000        3.500     0.000        0.000
dashboard_and_storytelling_skills           0.071        8.000          12.500        -4.500          0.570        3.500     0.000        0.000

dsj_rows  dsj_median_salary  dsj_median_min_exp  dsj_openings  an_postings  an_median_salary_mid  an_median_min_exp
seniority                                                                                                                      
Junior/Mid   803.000              9.300               1.000    68,797.000          NaN                   NaN                NaN
Senior       792.000             15.500               4.000    24,008.000    2,527.000                20.000             10.000
Junior           NaN                NaN                 NaN           NaN    7,509.000                 4.500              2.000
Mid              NaN                NaN                 NaN           NaN    4,713.000                12.500              5.000Share of postings mentioning each skill dimension:
                                  all_postings              core_data_roles             
seniority                               Junior   Mid Senior          Junior   Mid Senior
big_data_skills                          0.039 0.074  0.074           0.077 0.144  0.169
maths_stats_skills                       0.118 0.142  0.129           0.194 0.212  0.219
coding_skills                            0.217 0.294  0.207           0.280 0.373  0.345
ai_and_ml_skills                         0.067 0.091  0.077           0.144 0.183  0.170
dashboard_and_storytelling_skills        0.071 0.082  0.055           0.106 0.119  0.068                                   jds_cohens_d  jds_auc_single  demand_all  demand_junior  salary_premium_junior_LPA  effect_rank  demand_rank  rank_gap
big_data_skills                           0.223           0.561       0.056          0.039                      3.500        5.000        5.000     0.000
maths_stats_skills                        1.202           0.775       0.128          0.118                      3.500        2.000        2.000     0.000
coding_skills                             0.976           0.744       0.240          0.217                      3.500        3.000        1.000    -2.000
ai_and_ml_skills                          0.863           0.682       0.076          0.067                      3.500        4.000        4.000     0.000
dashboard_and_storytelling_skills         1.302           0.794       0.071          0.071                      3.500        1.000        3.000     2.000<Figure size 750x550 with 1 Axes>Caution: demand shares are keyword proxies on truncated text; treat the gap as a hypothesis to test, not a conclusion.

Soft-skill language in postings with a description (n=11,460); proxy mapping is an assumption:
                        sds_cohens_d  Junior   Mid  Senior
neuroticism                   -0.012   0.005 0.006   0.006
extraversion                   1.115   0.156 0.133   0.115
openness_to_experience         1.774   0.075 0.051   0.033
agreeableness                  0.595   0.014 0.012   0.013
conscientiousness              1.815   0.025 0.018   0.018
\end{lstlisting}

\subsubsection*{G.10 Dimensionality Reduction and Archetype Clustering}
\phantomsection\label{lst:clustering}
\textit{\hyperref[sec:multivariate]{[$\leftarrow$ Return to Archetype Clustering in Section V-D]}}

\begin{lstlisting}[style=pythonstyle,caption={PCA loadings, silhouette evaluation across k, and cluster centroid profiling.},label={lst:clustering_code}]
from sklearn.preprocessing import StandardScaler
from sklearn.decomposition import PCA
from sklearn.cluster import KMeans
from sklearn.metrics import silhouette_score

# JDS Clustering
X_j = StandardScaler().fit_transform(jds_core[JDS_SKILLS])
pca_j = PCA(n_components=3).fit(X_j)
sil_j = {k: silhouette_score(X_j, KMeans(k, random_state=42).fit_predict(X_j)) for k in range(2, 7)}
jds_core['cluster'] = KMeans(3, random_state=42).fit_predict(X_j)
prof_j = jds_core.groupby('cluster')[JDS_SKILLS + ['outcome']].mean()

# SDS Clustering
X_s = StandardScaler().fit_transform(sds_core[SDS_TRAITS])
pca_s = PCA(n_components=3).fit(X_s)
sil_s = {k: silhouette_score(X_s, KMeans(k, random_state=42).fit_predict(X_s)) for k in range(2, 7)}
sds_core['cluster'] = KMeans(4, random_state=42).fit_predict(X_s)
prof_s = sds_core.groupby('cluster')[SDS_TRAITS + ['success_classification_high_low']].mean()
\end{lstlisting}

\begin{lstlisting}[style=outputstyle,caption={Execution Output G.10: PCA Loadings and Cluster Profiles.},label={out:clustering}]
JDS: explained variance
     share
PC1  0.408
PC2  0.215
PC3  0.160
PC4  0.124
PC5  0.094JDS: loadings (PC1-PC3)
                            PC1    PC2    PC3
big data                 -0.026  0.903  0.237
maths stats               0.540 -0.206 -0.169
coding                    0.509  0.193 -0.550
ai & ml                   0.438 -0.207  0.779
dashboard & storytelling  0.506  0.251  0.073<Figure size 1200x450 with 2 Axes>JDS: silhouette by k -> {2: np.float64(0.333), 3: np.float64(0.337), 4: np.float64(0.283), 5: np.float64(0.263), 6: np.float64(0.315)} | chosen k=3
JDS: cluster profiles (pct_high = % in the high-outcome group)
         big_data_skills  maths_stats_skills  coding_skills  ai_and_ml_skills  dashboard_and_storytelling_skills   n  pct_high
cluster                                                                                                                       
0                  3.813               3.639          3.970             3.230                              3.691  23     4.348
1                  3.831               4.721          4.815             4.827                              4.877  78    83.333
2                  3.911               3.813          3.326             4.839                              3.687  38    18.421<Figure size 800x420 with 2 Axes>SDS: explained variance
     share
PC1  0.413
PC2  0.210
PC3  0.167
PC4  0.115
PC5  0.096SDS: loadings (PC1-PC3)
                          PC1    PC2    PC3
neuroticism            -0.247  0.766  0.483
extraversion            0.533 -0.040 -0.227
openness to experience  0.471  0.513 -0.229
agreeableness           0.374 -0.341  0.810
conscientiousness       0.541  0.178  0.083<Figure size 1200x450 with 2 Axes>SDS: silhouette by k -> {2: np.float64(0.313), 3: np.float64(0.322), 4: np.float64(0.335), 5: np.float64(0.212), 6: np.float64(0.216)} | chosen k=4
SDS: cluster profiles (pct_high = % in the high-outcome group)
         neuroticism  extraversion  openness_to_experience  agreeableness  conscientiousness   n  pct_high
cluster                                                                                                   
0             54.818        36.136                  42.182         34.545             35.682  22    27.273
1             35.950        36.500                  27.250         58.600             34.050  20     0.000
2             33.820        49.517                  47.955         48.225             53.674  89    88.764
3             29.733        34.133                  30.433         31.900             34.533  30     0.000<Figure size 800x460 with 2 Axes>
\end{lstlisting}

\subsubsection*{G.11 Repeated Cross-Validation and Interpretable Classifiers}
\phantomsection\label{lst:sds_models}\phantomsection\label{lst:jds_models}
\textit{\hyperref[sec:sds_models]{[$\leftarrow$ Return to SDS Models (V-A)]}} \quad \textit{\hyperref[sec:jds_models]{[$\leftarrow$ Return to JDS Models (V-B)]}}

\begin{lstlisting}[style=pythonstyle,caption={Repeated stratified 5-fold cross-validation and shallow decision tree rule extraction.},label={lst:models_code}]
from sklearn.model_selection import RepeatedStratifiedKFold, cross_validate
from sklearn.linear_model import LogisticRegression
from sklearn.tree import DecisionTreeClassifier, export_text
from sklearn.ensemble import RandomForestClassifier
from sklearn.inspection import permutation_importance

cv = RepeatedStratifiedKFold(n_splits=5, n_repeats=10, random_state=42)

# Models evaluation function
def evaluate_cv(X, y):
    models = {
        'Baseline': DummyClassifier(strategy='most_frequent'),
        'Logistic': LogisticRegression(),
        'Decision Tree': DecisionTreeClassifier(max_depth=3),
        'Random Forest': RandomForestClassifier(n_estimators=300, random_state=42)
    }
    return {name: cross_validate(m, X, y, cv=cv, scoring=['roc_auc', 'accuracy']) for name, m in models.items()}
\end{lstlisting}

\begin{lstlisting}[style=outputstyle,caption={Execution Output G.11a: Junior Model CV and Decision Rules (Table VII).},label={out:jds_models}]
JDS: predict high salary hike: repeated 5-fold CV (10 repeats)
                         cv_auc  auc_sd  cv_accuracy
model                                               
Baseline (majority)       0.500   0.000        0.525
Logistic (standardised)   0.903   0.053        0.857
Decision tree (depth 3)   0.800   0.082        0.782
Random forest             0.875   0.059        0.824JDS: logistic coefficients (per 1 SD)
                                   std_coef  odds_ratio_per_1sd
maths_stats_skills                    1.284               3.612
dashboard_and_storytelling_skills     1.117               3.057
ai_and_ml_skills                      0.762               2.143
big_data_skills                       0.683               1.980
coding_skills                         0.533               1.704JDS: random-forest permutation importance (30% hold-out)
                                   perm_importance    sd
maths_stats_skills                           0.091 0.038
dashboard_and_storytelling_skills            0.085 0.034
ai_and_ml_skills                             0.030 0.026
big_data_skills                              0.009 0.012
coding_skills                                0.005 0.018JDS: depth-3 tree rules
|--- dashboard_and_storytelling_skills <= 4.15
|   |--- maths_stats_skills <= 4.60
|   |   |--- class: 0
|   |--- maths_stats_skills >  4.60
|   |   |--- dashboard_and_storytelling_skills <= 2.55
|   |   |   |--- class: 0
|   |   |--- dashboard_and_storytelling_skills >  2.55
|   |   |   |--- class: 0
|--- dashboard_and_storytelling_skills >  4.15
|   |--- maths_stats_skills <= 3.55
|   |   |--- class: 0
|   |--- maths_stats_skills >  3.55
|   |   |--- coding_skills <= 4.65
|   |   |   |--- class: 1
|   |   |--- coding_skills >  4.65
|   |   |   |--- class: 1

JDS sensitivity: logistic AUC all rows=0.903 | excluding flagged rows (n=113)=0.875
\end{lstlisting}

\begin{lstlisting}[style=outputstyle,caption={Execution Output G.11b: Senior Model CV and Decision Rules (Table V).},label={out:sds_models}]
SDS: predict high success: repeated 5-fold CV (10 repeats)
                         cv_auc  auc_sd  cv_accuracy
model                                               
Baseline (majority)       0.500   0.000        0.528
Logistic (standardised)   0.964   0.033        0.927
Decision tree (depth 3)   0.942   0.039        0.934
Random forest             0.996   0.008        0.949SDS: logistic coefficients (per 1 SD)
                        std_coef  odds_ratio_per_1sd
conscientiousness          2.094               8.114
openness_to_experience     2.043               7.717
extraversion               0.954               2.595
neuroticism                0.798               2.222
agreeableness              0.628               1.873SDS: random-forest permutation importance (30% hold-out)
                        perm_importance    sd
openness_to_experience            0.081 0.035
conscientiousness                 0.044 0.021
agreeableness                     0.021 0.009
extraversion                      0.006 0.008
neuroticism                       0.003 0.002SDS: depth-3 tree rules
|--- openness_to_experience <= 38.50
|   |--- class: 0
|--- openness_to_experience >  38.50
|   |--- conscientiousness <= 36.50
|   |   |--- class: 0
|   |--- conscientiousness >  36.50
|   |   |--- agreeableness <= 37.50
|   |   |   |--- class: 0
|   |   |--- agreeableness >  37.50
|   |   |   |--- class: 1

SDS sensitivity: logistic AUC all rows=0.964 | excluding flagged rows (n=143)=0.955
\end{lstlisting}

\subsubsection*{G.12 Machine Learning Salary Regression Baselines}
\phantomsection\label{lst:salary_reg}
\textit{\hyperref[sec:salary_model]{[$\leftarrow$ Return to Salary Regressions in Section V-F]}}

\begin{lstlisting}[style=pythonstyle,caption={Cross-validated Ridge and HistGradientBoosting regressions with permutation importances.},label={lst:salary_reg_code}]
from sklearn.linear_model import Ridge
from sklearn.ensemble import HistGradientBoostingRegressor
from sklearn.model_selection import KFold

kf = KFold(n_splits=5, shuffle=True, random_state=42)

# DataScience Jobs Model
X_dsj = pd.get_dummies(dsj_core[['min_experience', 'role_family', 'num_of_jobs', 'is_senior_title']], drop_first=True)
y_dsj = np.log(dsj_core['avg_salary_lpa'])
ridge_dsj = cross_validate(Ridge(), X_dsj, y_dsj, cv=kf, scoring=['r2', 'neg_mean_absolute_error'])

# Analytics Jobs Model
X_an = pd.get_dummies(an_core[['exp_min_yrs', 'role_family', 'location_group'] + [c for c in an_core if c.startswith('skill_')]], drop_first=True)
y_an = an_core['salary_mid']
gbr_an = cross_validate(HistGradientBoostingRegressor(random_state=42), X_an, y_an, cv=kf, scoring=['r2', 'neg_mean_absolute_error'])
\end{lstlisting}

\begin{lstlisting}[style=outputstyle,caption={Execution Output G.12: Salary Regression Benchmarks and Feature Importance.},label={out:salary_reg}]
DataScience Jobs salary model (target=log avg_salary_lpa)
                   cv_r2  cv_mae
model                           
Baseline (mean)   -0.002   0.487
Ridge              0.571   0.305
Gradient boosting  0.552   0.311                 permutation_importance (R2 drop)
min_experience                              0.504
role_family                                 0.435
log_jobs                                    0.044
is_senior_title                             0.018
large_recruiter                             0.004Analytics Jobs salary model (target=salary_mid)
                   cv_r2  cv_mae
model                           
Baseline (mean)   -0.001   6.972
Ridge              0.439   5.022
Gradient boosting  0.444   4.989                                         permutation_importance (R2 drop)
exp_min_yrs                                                         0.798
role_family                                                         0.021
location_group                                                      0.014
skill_ai_and_ml_skills                                              0.004
n_skill_dims                                                        0.004
skill_maths_stats_skills                                            0.003
is_metro                                                            0.001
skill_coding_skills                                                 0.000
skill_dashboard_and_storytelling_skills                             0.000
skill_big_data_skills                                              -0.000Read these as descriptive drivers of advertised salary, not causal effects.
\end{lstlisting}

\subsubsection*{G.13 Export of Reproducible Artifacts}
\phantomsection\label{lst:export}
\textit{\hyperref[sec:pipeline]{[$\leftarrow$ Return to Reproducibility in Section IV]}}

\begin{lstlisting}[style=pythonstyle,caption={Automated serialization of clean tables, bridge matrices, and figures.},label={lst:export_code}]
exports = {
    'jds_core': jds_core, 'sds_core': sds_core,
    'dsj_core': dsj_core, 'analytics_core_no_text': an_core,
    'table_jds_effects': jds_effects, 'table_sds_effects': sds_effects,
    'table_dsj_senior_premium': prem, 'bridge_skill_alignment': alignment
}
for name, df in exports.items():
    df.to_csv(f'{name}.csv', index=True)
shutil.make_archive('dd_foundation_outputs', 'zip', OUT_DIR)
print('Saved all artifacts successfully.')
\end{lstlisting}

\begin{lstlisting}[style=outputstyle,caption={Execution Output G.13: Export Archive Generation Trace.},label={out:export}]
Saved 12 tables + 17 figures to outputs -> dd_foundation_outputs.zip
\end{lstlisting}

\subsubsection*{G.14 Final Analytical Recap and Verification}
\phantomsection\label{lst:recap}
\textit{\hyperref[sec:results]{[$\leftarrow$ Return to Analytical Results in Section VI]}}

\begin{lstlisting}[style=pythonstyle,caption={One-screen summary log confirming statistical signal integrity.},label={lst:recap_code}]
print('=' * 78)
print('FOUNDATION RECAP')
print('=' * 78)
print(f'Core rows -> JDS {len(jds_core)}, SDS {len(sds_core)}, DS Jobs {len(dsj_core)}, Analytics {len(an_core):,}')
print(f'Strongest junior skill signal : {jds_effects.loc[0, "feature"]} (d={jds_effects.loc[0, "cohens_d"]:.2f})')
print(f'Strongest senior trait signal : {sds_effects.loc[0, "feature"]} (d={sds_effects.loc[0, "cohens_d"]:.2f})')
print(f'Median senior salary premium across role families: {prem["senior_premium_pct"].median():.0f}%')
print(f'Largest importance-vs-demand gap: {alignment.sort_values("rank_gap", ascending=False).iloc[0]["feature"]}')
\end{lstlisting}

\begin{lstlisting}[style=outputstyle,caption={Execution Output G.14: Final Terminal Recap Log.},label={out:recap}]
==============================================================================
FOUNDATION RECAP
==============================================================================
Core rows -> JDS 139, SDS 161, DS Jobs 1595, Analytics 14,749
Strongest junior skill signal : dashboard & storytelling (d=1.30)
Strongest senior trait signal : conscientiousness (d=1.81); weakest: neuroticism (d=-0.01)
JDS model CV AUC (logistic) 0.90 | SDS 0.96
Median senior salary premium across role families: 64%
Largest importance-vs-demand gap: dashboard & storytelling (rank gap +2)
\end{lstlisting}

\subsection*{F. Acknowledgments}
The authors acknowledge the organizers of the Build for Bharat 2.0 Hackathon and the providers of the foundational workforce datasets.

\end{document}
